\usepackage[utf8]{inputenc}
\usepackage[T1]{fontenc}
\usepackage[spanish,es-tabla]{babel}
\usepackage{lmodern}
\usepackage[left=3cm,right=3cm,top=2.5cm,bottom=2.5cm,headheight=28pt]{geometry}
\usepackage{setspace}
\onehalfspacing
\usepackage{microtype}
\usepackage{graphicx}
\graphicspath{{imagenes/}}
\usepackage{amsmath,amssymb}
\usepackage{booktabs,longtable,tabularx}
\usepackage{float}
\usepackage{xcolor}
\usepackage{listings}
\usepackage{fancyhdr}
\usepackage{emptypage}
\usepackage[hidelinks,unicode]{hyperref}
\hypersetup{pdftitle={\thesistitle},pdfauthor={\thesisauthor}}
\usepackage[capitalise,noabbrev]{cleveref}
\crefname{chapter}{capítulo}{capítulos}
\crefname{section}{sección}{secciones}
\crefname{figure}{figura}{figuras}
\crefname{table}{tabla}{tablas}
\crefname{equation}{ecuación}{ecuaciones}
\crefname{lstlisting}{listado}{listados}
\renewcommand{\lstlistingname}{Listado}
\setcounter{secnumdepth}{3}
\setcounter{tocdepth}{2}
\pagestyle{fancy}
\fancyhf{}
\fancyhead[L]{\small\nouppercase{\leftmark}}
\fancyfoot[L]{\small\thesisauthor}
\fancyfoot[R]{\thepage}
\renewcommand{\headrulewidth}{0pt}
\fancypagestyle{plain}{\fancyhf{}\fancyfoot[L]{\small\thesisauthor}\fancyfoot[R]{\thepage}\renewcommand{\headrulewidth}{0pt}}
\lstset{basicstyle=\ttfamily\small,breaklines=true,numbers=left,numberstyle=\tiny,numbersep=8pt,columns=fullflexible,keepspaces=true,showstringspaces=false,captionpos=b,tabsize=2}
% Las notas son guías de redacción: eliminarlas al completar cada apartado.
\newcommand{\pendiente}[1]{\par\noindent\textit{Pendiente: #1}\par}
